% Generated draft: requires mathematical and source review.
\documentclass[11pt]{article}
\usepackage[margin=25mm]{geometry}
\usepackage{fontspec}
\setmainfont{Latin Modern Roman}
\usepackage{amsmath,amssymb,mathtools,unicode-math}
\setmathfont{Latin Modern Math}
\usepackage{enumitem,needspace,xurl,hyperref}
\setlength{\parindent}{0pt}
\setlength{\parskip}{5pt}
\setlength{\emergencystretch}{3em}
\newcommand{\N}{\mathbb{N}}
\newcommand{\Z}{\mathbb{Z}}
\newcommand{\Q}{\mathbb{Q}}
\newcommand{\R}{\mathbb{R}}
\newcommand{\C}{\mathbb{C}}
\newcommand{\F}{\mathbb{F}}
\newcommand{\A}{\mathbb{A}}
\newcommand{\PP}{\mathbb{P}}
\newcommand{\E}{\mathbb{E}}
\newcommand{\eps}{\varepsilon}
\newcommand{\dd}{\,\mathrm{d}}
\begin{document}
\section*{UnsolvedMath ID: 30002613}
\textbf{Source problem code:} catalogue entry
\subsection{Definitions and mathematical statement}
Let \(S\) be a projective K3 surface and let \(X=S^{[m]}\) denote its Hilbert scheme of length-\(m\) zero-dimensional subschemes. The source does not specify the base field or the range of \(m\). For codimension \(i\), write \(CH^i(X)\) for the Chow group of algebraic cycles modulo rational equivalence, using the source's unspecified coefficient convention, and write \(CH^i_{(j)}(X)\) for the components of the Chow decomposition intended by the source. Neither the construction of this decomposition nor its indexing convention is supplied. Let \(\operatorname{cl}_X\) denote the intended cycle-class map to \(H^{2i}(X)\); the cohomology theory and its coefficients are also unspecified.

The proposed assertion is that, for every projective K3 surface \(S\), every admissible \(m\), and every codimension \(0\leq i\leq 2m\), the intended decomposition satisfies
\[
CH^i_{(j)}(S^{[m]})=0\qquad\text{for every integer }j<0,
\]
and the restriction
\[
\operatorname{cl}_X\big|_{CH^i_{(0)}(X)}:
CH^i_{(0)}(X)\longrightarrow H^{2i}(X)
\]
is injective. Equivalently, a class in the grade-zero component with zero cycle class must be zero in the Chow group.

These assertions record the supplied question conditionally on the intended conventions. A fully specified conjecture requires identification of the decomposition, base field, Chow coefficients, cohomology theory and coefficients, and admissible values of \(m\); none is silently selected here.
\subsection{Short English statement}
For the Hilbert scheme of points on a projective K3 surface, the question asks whether all negative-grade Chow components vanish and whether every grade-zero cycle with zero cohomology class is rationally equivalent to zero, with the intended coefficients. The decomposition and several underlying conventions remain unspecified. The supplied record labels the problem partially solved; this draft does not independently verify that status.
\subsection{Sources}
\begin{itemize}
\item Supplied local UnsolvedMath registry record, permanent ID \(30002613\), titled \emph{Chow Decompositions of Hilbert Schemes of K3 Surfaces}; exact source problem code \texttt{catalogue entry}. The record belongs to the Oberwolfach Reports Open Problems collection and gives proposed year 2014. No individual proposer is identified.
\item Supplied report citation: \emph{Classical Algebraic Geometry (2015)}, Oberwolfach Reports, pp.\ 1695--1745, DOI \texttt{10.4171/owr/2014/31}: \url{https://doi.org/10.4171/owr/2014/31}. The collection describes these reports as documenting mathematical workshops and their open-problem discussions. The supplied title year 2015 and the 2014 DOI/proposed-year information are retained without reconciliation; no exact problem page is supplied.
\item Supplied literature reference: Maulik--Negut, \emph{Lehn's formula and Beauville--Voisin}, \url{https://arxiv.org/abs/1904.05262}.
\item Supplied literature reference: \emph{Beauville--Voisin--Franchetta conjecture and LLSS eightfolds}, \url{https://doi.org/10.1016/j.indag.2024.10.004}.
\item Supplied literature reference: Yin, \emph{Finite-dimensionality and K3 cycles}, \url{https://arxiv.org/abs/1404.0171}.
\item Literature assessment dated exactly \texttt{2026-08-21}: the supplied curation labels the problem partially solved and states that the general Hilbert-scheme-of-K3 case remains unresolved, while later work establishes weak or conditional forms and substantial bounded or motivic cases. It consequently describes the extracted 2014 claim as still accurate in general. This is supplied dated curation derived from the OWR record, not an independently verified literature assessment; the time-sensitive status and cited claims require independent verification.
\end{itemize}
\end{document}
